
\subsection{完全等方Log点を通る長鎖・非定数modular成分}
\label{subsec:isotropic-long-rankfour-membership}

前小節の階数四曲線を既存のrank一長鎖源へ直接同定する。
退化因子の一致だけで成分帰属を推測するのではなく、実際の
\(\mathrm{PGL}_5\)座標変換と全六Poisson括弧を与える。

\begin{theorem}[階数四成分への明示的な帰属]
\label{thm:isotropic-long-rankfour-membership}
定理\ref{thm:quadric-isotropic-dimension28}の
\(p=[P_0]\)と\(P_a=P_0+aK_1\)を用いる。
現行成分カタログの
\texttt{R4\_MODULAR\_RANK1\_LONG\_NONCONSTANT\_25}が定める
一般階数四、次元\(25\)の既約成分を
\(\mathcal Z_{\mathrm{long}}\)とする。このとき
\[
 [P_a]\in\mathcal Z_{\mathrm{long}}\quad(a\ne0),
 \qquad [P_0]\in\mathcal Z_{\mathrm{long}}
\]
が成り立つ。従って\(p\)は次元\(28\)の階数二成分
\texttt{R2\_LOG1112}と、この具体的に同定された
次元\(25\)の階数四成分との交点にある。
この同定には退化五次式の因子型だけでなく、
\(\mathrm{PGL}_5\)による全Poisson括弧の一致を用いる。

さらに、固定超平面\(u_1=0\)のaffine chartで
\(P_a\)のmodular場の射影斉次化の冪零Jordan型は、
\(a\ne0\)で\((3,1,1)\)、\(a=0\)で\((2,2,1)\)である。
\end{theorem}

\begin{proof}
同次座標を\((u_0,u_1,z,w_0,w_1)\)とし、
\(h=u_1\)のchartで
\[
 x=u_0/h,\qquad y=-w_0/h,\qquad
 z_{\mathrm{af}}=z/h,\qquad w=w_1/h
\]
と書く。以後\(z_{\mathrm{af}}\)を\(z\)と略記する。
定義\eqref{eq:isotropic-log-p0}と
\eqref{eq:isotropic-K}に代入すれば、
\(P_a=P_0+aK_1\)の非零括弧は厳密に
\begin{equation}\label{eq:isotropic-long-original-brackets}
\begin{aligned}
\{x,y\}&=a(y-x^2/2-zw),&
\{y,z\}&=-6z^2,\\
\{y,w\}&=2zw-2x^2+4y,&
\{z,w\}&=6z,
\end{aligned}
\qquad \{x,z\}=\{x,w\}=0.
\end{equation}
従って標準体積形式に関するmodular場は
\[
 M_a=-a\partial_x+(-ax+10z)\partial_y+10\partial_w.
\]
斉次座標ベクトルの基底を
\(e_h,e_x,e_y,e_z,e_w\)とすると、
このmodular場に対応する線形作用素\(N_a\)は
\[
 N_a(e_h)=-ae_x+10e_w,\qquad
 N_a(e_x)=-ae_y,\qquad
 N_a(e_z)=10e_y,\qquad N_a(e_y)=N_a(e_w)=0.
\]
\(a\ne0\)では\(\rank N_a=2,\ \rank N_a^2=1,\
N_a^3=0\)、\(a=0\)では\(\rank N_0=2,\ N_0^2=0\)。
従ってJordan型は述べた通りである。

\(a\ne0\)では次の\emph{アフィン線形}座標変換を施す：
\begin{equation}\label{eq:isotropic-long-projective-change}
 X=-x/a+10z/a^2,\quad Y=y/a^2,\quad
 Z=z,\quad W=w+10x/a.
\end{equation}
この変換のJacobianは\(-a^{-3}\ne0\)だから
射影線形自己同型へ延長する。とくに
\(y-zw\)のような\emph{非線形}座標変更を利用していない。

\(r=Y-X^2/2\)とし、
\[
 g=6Z,\qquad
 U=\frac{60Z}{a^2}\partial_W,\qquad d=\operatorname{div}U=0,
\quad
 K=-\frac{ZW}{a^2}+\frac{50Z^2}{a^4},\qquad
 F=(U+d+1)K=-\frac{ZW}{a^2}-\frac{10Z^2}{a^4}
\]
と置く。このとき\(U(g)=dg=0\)であり、
\(g\)は非定数一次式で、\(K\)は次数二以下である。
式\eqref{eq:isotropic-long-original-brackets}を
\eqref{eq:isotropic-long-projective-change}で変換した
\emph{全六括弧}は
\begin{equation}\label{eq:isotropic-long-source-brackets}
\begin{aligned}
 \{X,Y\}&=-F-(1+d)r,&
 \{X,Z\}&=U_Z=0,&
 \{X,W\}&=U_W=\frac{60Z}{a^2},\\
 \{Y,Z\}&=XU_Z+r\,g_W+(gK)_W=-\frac{6Z^2}{a^2},\\
 \{Y,W\}&=XU_W-r\,g_Z-(gK)_Z
   =\frac{60XZ}{a^2}-6r+\frac{12ZW}{a^2}
      -\frac{900Z^2}{a^4},\\
 \{Z,W\}&=g=6Z .
\end{aligned}
\end{equation}
これは既存の
\texttt{RANK\_ONE\_LONG\_NONCONSTANT\_COMPONENT\_ENTRANCE.md}
第1節にある全源の式(1)と、一つの係数も省略せず一致する。

実際、そのsourceのPfaffianは
\[
 \operatorname{Pf}(P_a)=-g(r+K)
   =-6Z\left(Y-\frac{X^2}{2}-\frac{ZW}{a^2}
                           +\frac{50Z^2}{a^4}\right).
\]
二次因子の斉次化
\[
 Q_a=hY-\frac{X^2}{2}-\frac{ZW}{a^2}
                              +\frac{50Z^2}{a^4}
\]
の半Hessianの行列式は\(-1/(32a^4)\ne0\)なので
一般階数四源の滑らかな二次因子という開条件も満たす。
その二重超平面\(h=0\)は\(Q_a=0\)に接し、
単純超平面\(Z=0\)は接点を通る。
従って全\(a\ne0\)に対して、実際の長鎖非定数\(g\)源の
\(\mathrm{PGL}_5\)飽和に属する。

既存のカタログでこの源の既約閉包
\(\mathcal Z_{\mathrm{long}}\)が次元\(25\)の
Poisson既約成分であるという結果を用いれば
\([P_a]\in\mathcal Z_{\mathrm{long}}\)となる。
元の同次座標で\(P_a=P_0+aK_1\)は\(a=0\)まで
多項式的なPoisson曲線であり、\(\mathcal Z_{\mathrm{long}}\)
は閉だから\([P_0]\)もこれに含まれる。
階数二対数型成分への帰属とその次元\(28\)は
定理\ref{thm:quadric-isotropic-dimension28}で既に証明した。
\end{proof}

\begin{remark}[既知の成分性と今回の同定の分離]
\texttt{R4\_MODULAR\_RANK1\_LONG\_NONCONSTANT\_25}が
既約成分であるという入力は、本リポジトリの先行する
全源成分証明
\texttt{research/classification/RANK\_ONE\_LONG\_NONCONSTANT\_25\_COMPONENT.md}
に依存する。本節が新たに証明するのは
\(P_0+aK_1\)のこの\emph{特定の}源への明示同定である。
\(p\)を通る他の階数四成分の不存在、
全接錐の分解、ambient局所環の被約性は示していない。
\end{remark}


\begin{corollary}[Poissonスキームの非等次元性と非Cohen--Macaulay性]
\label{cor:poisson-p4-not-cm-at-isotropic-log}
定理\ref{thm:quadric-isotropic-dimension28}と
定理\ref{thm:isotropic-long-rankfour-membership}の点
\(p=[P_0]\in\mathcal P\)について、
\(\mathcal O_{\mathcal P,p}\)は非等次元であり、
\[
 \dim\mathcal O_{\mathcal P,p}=28,\qquad
 \operatorname{depth}\mathcal O_{\mathcal P,p}\le25.
\]
従ってCohen--Macaulay欠損
\(\dim-\operatorname{depth}\)は少なくとも\(3\)であり、
Gorensteinでも局所完全交叉でもない。
被約化した局所環についても同じ深さの上界が成立する。
\end{corollary}
\begin{proof}
階数二対数型既約成分\(\mathcal Z_{\mathrm{log}}\)は
次元\(28\)、長鎖非定数型の
\(\mathcal Z_{\mathrm{long}}\)は次元\(25\)であり、
両者は\(p\)を含む相異なる\emph{既約成分}である。
\(p\)は複素射影スキームの閉点なので、
それぞれの局所環の次元も\(28\)、\(25\)であり、
\(\mathcal O_{\mathcal P,p}\)の極小素イデアルのうち
少なくとも二つは異なる次元の商を与える。
従って局所環は非等次元である。
この\(25\)次元成分に対応する極小素イデアル
\(\mathfrak p_{\mathrm{long}}\)は
Noether局所環のassociated primeでもある。
associated primeに対する一般の深さ不等式
\(\operatorname{depth}R\le\dim(R/\mathfrak p)\)
（\cite{StacksDepth}, Lemma 10.72.9）より
\(\operatorname{depth}\mathcal O_{\mathcal P,p}\le25\)。
局所次元は\(28\)であるから、
Cohen--Macaulay欠損は少なくとも\(3\)である。
被約化しても二つの極小素イデアルと商次元は変わらず、
同じ深さ上界が成立する。
Gorenstein局所環と局所完全交叉環は
Cohen--Macaulayであるので、これらでもない。
\end{proof}

\begin{remark}
非等次元性は被約台にも見られる。
従ってこの系はambient Poissonスキームの
非被約性を仮定も証明もしない。
二つ以外の既約成分が\(p\)を通るか、
全接錐の既約分解がどうなるかは未解決である。
\end{remark}
